% ============================================================================
%  CONTRATO DE NOTACION — Aprendizaje Automatico: Prediccion (G244)
%  Se incluye en TODOS los capitulos. No redefinir simbolos localmente.
%
%  Solo sobreviven al  dentro de $$...$$:
%      \newcommand, \renewcommand, \DeclareMathOperator
%  NO: \usepackage, \DeclarePairedDelimiter, \newenvironment.
% ============================================================================

% ---- EL INTERRUPTOR ---------------------------------------------------------
% Todo coeficiente pasa por \coef. Cambiar w por \beta es UNA linea, no 15
% ficheros. Se usa \boldsymbol (no \mathbf) para que funcione con griegas.
\newcommand{\coef}{\boldsymbol{w}}
\newcommand{\coefhat}{\hat{\boldsymbol{w}}}
\newcommand{\coefj}{w_j}
\newcommand{\coefzero}{w_0}
\newcommand{\coefols}{\hat{\boldsymbol{w}}_{\mathrm{OLS}}}
\newcommand{\coefridge}{\hat{\boldsymbol{w}}_{\lambda}}

% ---- Datos y matriz de diseno ----------------------------------------------
\newcommand{\X}{\mathbf{X}}
\newcommand{\Xt}{\mathbf{X}^{\mathsf{T}}}
\newcommand{\xv}{\mathbf{x}}
\newcommand{\yv}{\mathbf{y}}
\newcommand{\Phimat}{\boldsymbol{\Phi}}
\newcommand{\onesv}{\mathbf{1}}
\newcommand{\zerov}{\mathbf{0}}
\newcommand{\Id}{\mathbf{I}}
\newcommand{\Zbf}{\mathbf{Z}}
\newcommand{\Vbf}{\mathbf{V}}
\newcommand{\Hess}{\nabla^2}
\newcommand{\Hbf}{\mathbf{H}}
\newcommand{\Wbf}{\mathbf{W}}

% ---- Targets, residuos ------------------------------------------------------
\newcommand{\yi}{y_i}
\newcommand{\yhat}{\hat{y}}
\newcommand{\yhati}{\hat{y}_i}
\newcommand{\ybar}{\bar{y}}
\newcommand{\resid}{\mathbf{r}}
\newcommand{\residi}{r_i}

% ---- Tamanos ----------------------------------------------------------------
\newcommand{\nobs}{n}
\newcommand{\nfeat}{p}
\newcommand{\ntrain}{n_{\mathrm{train}}}
\newcommand{\nval}{n_{\mathrm{val}}}
\newcommand{\ntest}{n_{\mathrm{test}}}
\newcommand{\ncal}{n_{\mathrm{cal}}}

% ---- Riesgo y perdidas ------------------------------------------------------
% Los acronimos van en INGLES (MSE, RMSE, MAE, RSS, SST...): es lo que veran en
% scikit-learn y en la bibliografia. El nombre completo, en la prosa, va en espanol.
% \SSR es SS_reg y no SSR, porque en ingles SSR se usa tambien para la suma de
% cuadrados de los residuos, que aqui es \RSS.
% R = riesgo VERDADERO ; \Riskh = riesgo EMPIRICO. La distincion es del cap. 8.
\newcommand{\Risk}{R}
\newcommand{\Riskh}{\hat{R}}
\newcommand{\Riskproc}[1]{\bar{R}_{#1}}
\newcommand{\loss}{\ell}
\newcommand{\MSE}{\mathrm{MSE}}
\newcommand{\RMSE}{\mathrm{RMSE}}
\newcommand{\MAE}{\mathrm{MAE}}
\newcommand{\RSS}{\mathrm{RSS}}
\newcommand{\SSE}{\mathrm{SSE}}
\newcommand{\SST}{\mathrm{SST}}
\newcommand{\SSR}{\mathrm{SS_{reg}}}
\newcommand{\Rsq}{R^2}
\newcommand{\CV}{\mathrm{CV}}

% ---- Verosimilitud ----------------------------------------------------------
\newcommand{\Lik}{L}
\newcommand{\loglik}{\ell}
\newcommand{\logliksp}{\ell_{\mathrm{p}}}
\newcommand{\pdf}{p}
\newcommand{\Normal}{\mathcal{N}}
\newcommand{\Laplacedist}{\mathrm{Laplace}}
\newcommand{\given}{\,\vert\,}
\newcommand{\iid}{\overset{\text{iid}}{\sim}}

% ---- Escalares fijos --------------------------------------------------------
\newcommand{\noise}{\sigma}
\newcommand{\noisehat}{\hat{\sigma}}
\newcommand{\lr}{\alpha}          % tasa de aprendizaje
\newcommand{\reg}{\lambda}        % OJO: en scikit-learn esto se llama alpha=
\newcommand{\defeq}{\coloneqq}

% ---- Calculo ----------------------------------------------------------------
\newcommand{\grad}{\nabla}
\newcommand{\pd}[2]{\frac{\partial #1}{\partial #2}}
\newcommand{\dd}{\mathrm{d}}
\DeclareMathOperator*{\argmax}{arg\,max}
\DeclareMathOperator*{\argmin}{arg\,min}

% ---- Algebra ----------------------------------------------------------------
\newcommand{\T}{^{\mathsf{T}}}
\newcommand{\norm}[1]{\left\lVert #1 \right\rVert}
\newcommand{\abs}[1]{\left\lvert #1 \right\rvert}
\newcommand{\inner}[2]{\left\langle #1, #2 \right\rangle}
\DeclareMathOperator{\rank}{rank}
\DeclareMathOperator{\tr}{tr}
\DeclareMathOperator{\diag}{diag}
\DeclareMathOperator{\kernelop}{null}
\DeclareMathOperator{\sign}{sign}
\newcommand{\Real}{\mathbb{R}}
\newcommand{\Realn}{\mathbb{R}^{n}}
\newcommand{\Realp}{\mathbb{R}^{p+1}}
\newcommand{\Rbb}{\mathbb{R}}

% ---- Probabilidad -----------------------------------------------------------
\newcommand{\E}[1]{\mathbb{E}\!\left[ #1 \right]}
\newcommand{\Esub}[2]{\mathbb{E}_{#1}\!\left[ #2 \right]}
\newcommand{\Var}[1]{\mathrm{Var}\!\left( #1 \right)}
\newcommand{\Varsub}[2]{\mathrm{Var}_{#1}\!\left( #2 \right)}
\newcommand{\Cov}[2]{\mathrm{Cov}\!\left( #1, #2 \right)}
\newcommand{\Prob}[1]{\mathbb{P}\!\left( #1 \right)}
\DeclareMathOperator{\med}{median}

% ---- Conjuntos y modelos ----------------------------------------------------
\newcommand{\data}{\mathcal{D}}
\newcommand{\Dtrain}{\mathcal{D}_{\mathrm{train}}}
\newcommand{\Dval}{\mathcal{D}_{\mathrm{val}}}
\newcommand{\Dajuste}{\mathcal{D}_{\mathrm{ajuste}}}
\newcommand{\Dtest}{\mathcal{D}_{\mathrm{test}}}
\newcommand{\fold}[1]{\mathcal{F}_{#1}}
\newcommand{\model}{f}
\newcommand{\modelh}{\hat{f}}
\newcommand{\nullmodel}{\bar{f}}
\newcommand{\hyp}{\mathcal{F}}
\newcommand{\Xcal}{\mathcal{X}}
\newcommand{\Ycal}{\mathcal{Y}}

% ---- Adelantos (definidos ya, usados a partir del cap. 8) -------------------
\newcommand{\pos}[1]{\left( #1 \right)_{+}}
\newcommand{\softthr}{S}
\newcommand{\basis}{\varphi}
% Mapa de caracteristicas del cap. 6: componente escalar, vector y numero.
\newcommand{\basisv}{\boldsymbol{\varphi}}
\newcommand{\nbasis}{d}
\newcommand{\pinball}{\rho_{\tau}}
